\chapter{MARCO TEÓRICO}

Este capítulo tiene dos partes. La primera presenta los fundamentos sobre los que se
construye CAVA: la segmentación por desplazamiento de embeddings, el consenso de particiones
mediante una matriz de co-asociación, el agrupamiento jerárquico y las métricas de
evaluación. La segunda describe el diseño del sistema por módulos: la extracción de señales,
el gestor multimodal sin fusión, el mapa de tópicos, la voz y los hablantes, y las versiones
por consumidor. La Figura~\ref{fig:mapa-mental} muestra el flujo completo: las
señales se calculan en paralelo, cada una vota y ninguna representación se combina con otra;
el mapa de tópicos sale del consenso y alimenta los formatos de consumidor.

\begin{figure}[H]
\centering
\includegraphics[width=\textwidth]{generadas/f21_flujo_cava}
\caption{Flujo de CAVA. El video se divide en unidades de 5\,s; cada señal calcula sus
propias fronteras y etiquetas dentro de la franja roja, donde ninguna representación se
combina con otra; la votación decide fronteras y tópicos con quórum y abstención; del mapa de
tópicos salen los cinco formatos de versión.}
\label{fig:mapa-mental}
\end{figure}

\section{Fundamentos}

Los conceptos de este capítulo se ilustran con un caso sintético: un documental de cuatro
minutos con seis escenas de 40\,s, cada una de 8 unidades. Las dos primeras tratan de leones
(tema A), las dos siguientes del océano (tema B), la quinta es un paisaje sin narración (tema
C) y la última vuelve a los leones, de modo que el tema A se retoma. Las señales visual, de
diálogo y del VLM se generan a partir de un prototipo por tema más ruido, y el diálogo no
tiene evidencia durante la escena sin narración. Los fotogramas son ilustraciones creadas
para este ejemplo y no provienen de ningún corpus; el procesamiento, en cambio, es el del
gestor de CAVA con los parámetros del perfil de documental.

\subsection{Segmentación por desplazamiento de embeddings}
Sea $x_1, \dots, x_n$ la secuencia de embeddings de una señal, uno por unidad temporal. El
score de frontera entre las unidades $t$ y $t+1$ compara la media de una ventana de $w$
unidades a cada lado:
\begin{equation}
  b_t = 1 - \cos\!\Big(\tfrac{1}{w}\textstyle\sum_{k=0}^{w-1} x_{t-k},\;
        \tfrac{1}{w}\textstyle\sum_{k=1}^{w} x_{t+k}\Big),
\end{equation}
con los $x_t$ normalizados a norma unitaria y $w = 2$ en el gestor. El score se divide por su
máximo en el video. Los picos de $b_t$ son candidatos a frontera, como en TextTiling
\cite{hearst1997texttiling}.
CAVA corta donde $b_t$ supera un percentil propio de cada señal y exige una separación mínima
entre cortes. Usar un percentil y no un umbral absoluto hace que la cantidad de cortes no
dependa de la escala de cada embedding.
La Figura~\ref{fig:m1} muestra el score de cada señal en el caso ilustrativo.

\begin{figure}[H]
\centering
\includegraphics[width=\textwidth]{generadas/m1_fronteras_por_senal}
\caption{Score de frontera $b_t$ de cada señal en el caso ilustrativo, con el color del tema
de cada escena. Las señales visual y del VLM marcan sus picos en los cambios de tema (1:20,
2:40 y 3:20). El diálogo solo tiene un pico claro en 1:20; durante la escena sin narración
(gris) su score es nulo, y el percentil lo obliga a proponer cortes de poca altura en 2:35 y
3:55.}
\label{fig:m1}
\end{figure}

\subsection{Consenso de particiones y matriz de co-asociación}
Dadas $S$ particiones de las mismas $n$ unidades, cada una con etiquetas $\ell^{s}_i$, la
matriz de co-asociación cuenta en cuántas particiones aparece cada par en el mismo grupo
\cite{fred2005evidence}:
\begin{equation}
  V_{ij} = \sum_{s=1}^{S} w_s\, \mathbf{1}\big[\ell^{s}_i = \ell^{s}_j\big].
\end{equation}
Con $w_s = 1$, $V_{ij}/S$ es la fracción de particiones que ponen juntos a $i$ y $j$. La
partición de consenso se obtiene agrupando la matriz de distancias $D = 1 - V/S$. Este
esquema combina decisiones y no representaciones: solo necesita las etiquetas de cada
partición, nunca los embeddings que las produjeron \cite{strehl2002cluster}.

Cuando una partición no tiene información sobre una unidad (por ejemplo, la transcripción en
un tramo sin diálogo), contarla como ``separa a $i$ de $j$'' introduce un desacuerdo falso.
CAVA define el conjunto de votantes activos en un par, $A_{ij} = \sum_s w_s\,
\mathbf{1}[\ell^s_i \neq -1 \wedge \ell^s_j \neq -1]$, donde $-1$ indica abstención, y usa
la fracción $V_{ij}/A_{ij}$ en lugar de $V_{ij}/S$ (Figura~\ref{fig:abstencion}).

\begin{figure}[H]
\centering
\includegraphics[width=\textwidth]{generadas/m3_matriz_votos}
\caption{Matriz de votos del caso ilustrativo, ordenada en el tiempo (48 unidades). Los
bloques oscuros de la diagonal son los temas, y los bloques fuera de la diagonal entre las
dos apariciones de A son el retorno. El recuadro rojo marca la escena sin narración: sin
abstención, el diálogo vota sobre esas unidades como si su silencio fuera información; con
abstención, la fracción se calcula solo sobre la señal visual y el VLM.}
\label{fig:abstencion}
\end{figure}

\subsection{Agrupamiento jerárquico y percolación}
El agrupamiento jerárquico aglomerativo une en cada paso los dos grupos más cercanos. Con
enlace simple, la distancia entre grupos es la del par más cercano; con enlace promedio, la
media de todos los pares. Si se umbraliza $V$ y se toman las componentes conexas del grafo
resultante, el resultado equivale al enlace simple: basta una cadena $i \sim j \sim k$ de
pares que superan el quórum para que $i$ y $k$ queden en el mismo grupo aunque casi ninguna
señal los agrupe. En una secuencia larga con retornos, estas cadenas se propagan hasta formar
un tópico gigante, fenómeno que aquí se llama \emph{percolación}. El enlace promedio exige
que el consenso se mantenga en promedio sobre todos los pares de los dos grupos, lo que corta
las cadenas \cite{fred2005evidence} (Figura~\ref{fig:m4}).

\begin{figure}[H]
\centering
\includegraphics[width=\textwidth]{generadas/m4_percolacion}
\caption{Percolación en una cadena de seis segmentos. Cada par vecino tiene el quórum de dos
señales, pero los extremos no tienen nada en común. Con componentes conexas la cadena forma
un solo tópico; con enlace promedio se separa en tres grupos, porque el consenso medio entre
grupos no alcanza el quórum.}
\label{fig:m4}
\end{figure}

\subsection{Métricas}
El coeficiente de correlación de Matthews \cite{matthews1975comparison} resume una matriz de
confusión binaria en un valor entre $-1$ y $1$:
\begin{equation}
  \mathrm{MCC} = \frac{TP\cdot TN - FP\cdot FN}
  {\sqrt{(TP+FP)(TP+FN)(TN+FP)(TN+FN)}}.
\end{equation}
CAVA lo aplica a la clasificación de pares de unidades (mismo tópico o no). A diferencia de la
exactitud o el F1, el MCC vale cero para un predictor que une todo o que separa todo, de modo
que no premia ni la sobresegmentación ni la percolación. El índice de Rand ajustado (ARI)
\cite{hubert1985comparing} compara dos particiones corrigiendo el acuerdo esperado por azar, y
se usa para la identidad de hablante. Para elegir el número de hablantes se usa el criterio de
información bayesiano (BIC) \cite{schwarz1978estimating} sobre un agrupamiento de Ward
\cite{ward1963hierarchical}, que en cada paso une los grupos cuya fusión menos aumenta la
varianza interna.

\section{Diseño del sistema: definiciones}

\begin{itemize}
  \item Unidad: tramo temporal fijo (5\,s en los perfiles de película y documental, 10\,s en
        reunión) sobre el que cada señal produce un embedding o un score. Es el elemento que
        se vota. Las fronteras finales se ajustan al corte visual certificado por SegLoop más
        cercano, dentro de una ventana de hasta 30\,s.
  \item Señal: fuente de evidencia que se procesa por separado hasta la decisión. El producto
        usa cuatro señales de embeddings (transcripción proyectada con bge, visual SigLIP
        \cite{zhai2023siglip} por unidad, voz con log-mel y tono, y descripciones de un VLM
        \cite{bai2025qwen25vl} por bloques de $\sim$40\,s) y, en el perfil de reunión, los
        turnos de hablante como votante de fronteras.
  \item Fusión: combinar embeddings o scores de señales distintas en una representación única
        antes de decidir. CAVA no la usa en ningún punto. La base de comparación
        Fusión de los experimentos concatena los embeddings normalizados de las
        señales y los pasa por el mismo gestor como si fueran una sola señal.
  \item Consenso de votos: cada señal emite sus propios cortes (percentil propio de su score
        de frontera) y sus propias etiquetas de tópico. Una frontera queda confirmada si
        alcanza la mayoría ponderada de las señales o un número mínimo fijo de señales;
        un par de unidades se enlaza si un quórum de señales lo
        agrupa.
  \item Tópico: grupo de unidades, posiblemente no contiguas, que un quórum de señales
        agrupa. Un retorno es la reaparición de un tópico después de una ausencia.
  \item Abstención: una señal sin evidencia en una unidad no vota sobre ella (etiqueta $-1$).
  \item Procedencia: por cada grupo, el aporte de cada señal medido como acuerdo (fracción de
        pares del grupo que la señal pone en su mismo tópico propio), soporte (fracción de
        unidades que comparten etiqueta con otra del grupo) y coherencia de sus embeddings
        dejando una unidad fuera.
  \item Versión por consumidor: producto derivado del mapa con un presupuesto declarado:
        capítulos, recapitulación de 60\,s, destacados del 10\,\%, hilos de tópico e hilos
        por hablante.
\end{itemize}

\section{Extracción de señales}

La transcripción usa faster-whisper large-v3
\cite{radford2023whisper} y se proyecta a embeddings con bge \cite{xiao2023bge}. La señal
visual codifica un frame por unidad con SigLIP. La voz extrae descriptores de log-mel y tono.
El VLM (Qwen2.5-VL-7B, congelado, sin ajuste) describe bloques de $\sim$40\,s a partir de
tres frames y del habla del bloque, y esas descripciones, proyectadas con bge, forman una
cuarta señal de embeddings. El diseño admite además otras señales con la misma interfaz;
el texto en pantalla (OCR) y el espectro de audio están definidos pero no implementados. Los perfiles de contenido fijan
el percentil de corte por dominio: p80 para película, p90 para documental y p95 para reunión.

\section{Gestor multimodal sin fusión}

El gestor, núcleo del sistema, ejecuta cuatro pasos:

\begin{enumerate}
  \item Por señal: score de frontera propio (ecuación del desplazamiento con ventana $w$),
        cortes propios por percentil con separación mínima y etiquetas de tópico propias
        sobre esos cortes.
  \item Consenso de fronteras: cada señal marca presencia en $\pm 1$ unidad alrededor de
        cada corte propio, los votos se suman con el peso de cada señal y cada racha de
        posiciones contiguas que alcanza la mayoría se reduce a un único corte en el máximo
        de votos. En el producto todas las señales de embeddings pesan 1.0; el votante de
        densidad de habla, usado en los pilotos, pesa 0.5.
  \item Consenso de enlace: la matriz de votos $V_{ij}$ se compara con el quórum $q$ y los
        grupos se forman como componentes conexas o por enlace promedio. Con abstención, un par solo se enlaza si tiene al
        menos $q$ votantes activos.
  \item Procedencia y banderas: cada grupo reporta el aporte de cada señal; un grupo de una
        unidad sin soporte de varias señales se marca como sin tópico, y un grupo con más
        segmentos que un máximo fijado se marca para revisión.
\end{enumerate}

En el caso ilustrativo (Figura~\ref{fig:m2}), las tres señales coinciden en 1:20, 2:40 y
3:20. El corte del diálogo en 2:35 cae dentro de la tolerancia del corte de 2:40 y se suma a
él, mientras que el de 3:55 queda con un solo voto y se descarta.

\begin{figure}[H]
\centering
\includegraphics[width=\textwidth]{generadas/m2_consenso_fronteras}
\caption{Consenso de fronteras en el caso ilustrativo. Arriba, los cortes propios de cada
señal y su ventana de tolerancia; abajo, los votos por posición. Solo las posiciones que
alcanzan la mayoría (dos de tres señales) producen un corte, ubicado en el máximo de votos de
cada racha.}
\label{fig:m2}
\end{figure}

\section{Retornos de tópico}

Las etiquetas de cada señal admiten retornos. En orden
temporal, cada segmento puede unirse al tópico anterior con el que tenga mayor similitud
coseno, si esa similitud supera un percentil de las similitudes entre segmentos no adyacentes
de la misma señal. Un segmento nunca vuelve al tópico del que acaba de separarse y solo puede
unirse a un tópico existente. De este modo, que un tema se retome es una decisión de cada
señal, que el quórum después confirma o descarta.
La Figura~\ref{fig:m5} muestra el resultado completo sobre el caso ilustrativo: las tres
señales reconocen el regreso a los leones, el mapa final une los dos tramos del tema A en un
solo tópico y la procedencia indica que el grupo del paisaje está sostenido solo por la señal
visual y el VLM, porque el diálogo se abstiene.

\begin{figure}[H]
\centering
\includegraphics[width=\textwidth]{generadas/m5_mapa_topicos}
\caption{Mapa de tópicos del caso ilustrativo. Las tres primeras filas son las etiquetas
propias de cada señal, coloreadas según el tema de sus unidades (rayado: abstención). La
fila \emph{mapa} es el resultado del gestor con enlace promedio y abstención, y el arco une
los dos tramos del tema A. Abajo, el soporte de cada señal en cada grupo.}
\label{fig:m5}
\end{figure}

\section{Mapa de tópicos}

El mapa de tópicos es la salida que usan los formatos. Por video tiene dos niveles de
granularidad (uno fino, con quórum de todas las señales, y uno por defecto, con quórum 2) y, por tópico, las unidades que lo forman, sus tramos en segundos, la duración
acumulada, el número de retornos, una etiqueta generada por un modelo de lenguaje
(Qwen2.5-14B) a partir de los tramos, la procedencia por señal, las banderas y el contexto de
cada tramo (la descripción del VLM y el diálogo, ambos producidos por el propio sistema). Los
nombres de personajes se propagan por los enlaces del mapa. Las identidades por apariencia
(P\#) se calculan pero no se publican como identidad.

\section{Agregador adaptativo}

Como alternativa al percentil fijo, el agregador adaptativo deja que la longitud de los
tópicos dependa de dónde las señales siguen viendo continuidad. Cada votante juzga cada par
de unidades vecinas comparando la similitud cruzada con su coherencia interna y un umbral
propio, la unión del par requiere mayoría estricta de votantes activos y cada ronda se
registra como un nivel de una jerarquía. Esta variante trata la granularidad como un
resultado y no como un hiperparámetro; no forma parte del producto congelado.

\section{Voz e identidad de hablante}

La voz entra al sistema de tres maneras. Como votante de tópico, sus embeddings de log-mel y
tono agrupan unidades acústicamente parecidas. Como votante de fronteras, la caída de la
densidad de habla marca pausas. Como identidad, los embeddings de hablante WavLM-sv
\cite{chen2022wavlm} se agrupan con Ward, el número de grupos se elige con BIC y el resultado
son las identidades V\#. Un detector de hablante único compara la dispersión de los
embeddings con un umbral $\tau$ calibrado y decide si el video tiene un solo hablante, como un
narrador. En el perfil de documental las identidades de voz se publican como grupos acústicos
y no como personas, porque el detector no reconoce al narrador único.

\section{Versiones por consumidor}

Del mapa se derivan cinco formatos con presupuestos verificables: capítulos (fronteras del
nivel superior), recapitulación de 60\,s (una unidad por tópico hasta llenar el presupuesto),
destacados del 10\,\% (tramos representativos de los tópicos más largos), hilos de tópico
(una lista de reproducción por tópico con sus retornos) e hilos por hablante (una lista por
identidad de voz). La compuerta G-versión verifica, por video y por formato, que cada lista
exista, cumpla su presupuesto y que cada perfil tenga al menos una versión. No mide la
calidad percibida, que corresponde a la evaluación con personas (EXP-040).
